\chapter{The Knowledge Environment}\label{ch:environment}

\section{Design Goals}\label{sec:env-goals}
\subsection{Decoupling the Graph Source from the Question Sets}
\subsection{Why a Curated Subgraph Rather Than LLM-Extracted Triples}

\section{Source Data}\label{sec:source-data}
\subsection{The Freebase Snapshot}
\subsection{WebQSP and ComplexWebQuestions as Question Sets}

\section{Graph Construction}\label{sec:graph-construction}
% NOTE: the graph is the UNION of RoG's pre-extracted per-question subgraphs
% (name-keyed triples from the WebQSP and CWQ parquet dumps), NOT a topic-entity
% seeding plus k-hop expansion performed here. This determines the coverage
% ceiling and is exactly what an examiner probes.
\subsection{Union of Per-Question Subgraphs from the RoG Distribution}
\subsection{Relation Filtering and Noise Removal}
\subsection{Treatment of Mediator Nodes and Its Effect on Hop Counts}

\section{Graph Store Construction}\label{sec:graph-store}
\subsection{Property-Graph Schema}
\subsection{Bulk Import into Neo4j}
\subsection{Graph Statistics}

\section{The Vector Index}\label{sec:vector-index}
\subsection{Entity Embeddings}
\subsection{Relation-Vocabulary Embeddings}
\subsection{Scope: Linking and Candidate Ranking, Never Ground Truth}

\section{Entity Resolution and Surface-Form Linking}\label{sec:entity-resolution}
% Call this cascade exact / lexical / vector. Never "Tier 1/2/3" --- "tier" is
% reserved for the groundedness metric in Chapter~\ref{ch:setup}.
\subsection{The Exact, Lexical, and Vector Cascade}
\subsection{Threshold Selection on Development Data}

\section{Environment Validation Gate}\label{sec:validation-gate}
\subsection{Answer-Reachability Protocol}
\subsection{Coverage Results and the Induced Accuracy Ceiling}
\subsection{Analysis of Linking Misses and Gate Failures}
\subsection{What the Environment Cannot Express}\label{sec:unanswerable}
% No date literals, no numeric literals, no ordinal encodings --- all three
% verified. Define the "unanswerable-in-environment" question class here;
% Section~\ref{sec:kg-gaps} invokes it rather than introducing it.

\endinput
